\centerline{\bf \Huge Банки. База(А+Д). КР}
\centerline{\bf \Huge }

\begin{flushright}
        \scriptsize{\textbf{tt}}
\end{flushright}


\problem{\textbf{1}} В июле некоторого года планируется взять кредит в банке на некоторую сумму. Условия его возврата таковы:

— каждый январь долг возрастает на 25\% по сравнению с долгом на конец предыдущего
года;

— с февраля по июнь необходимо выплатить часть долга одним платежом.
Известно, что сумма всех выплат составила 375 000 рублей. 

Сколько рублей было взято в банке, если известно, что кредит был полностью погашен четырьмя равными платежами?

\problem{\textbf{2}} 15 января планируется взять кредит в банке на 25 месяцев. Условия его возврата таковы:

— 1-ого числа каждого месяца долг возрастает на $r$\% по сравнению с концом предыдущего
месяца;

— со 2-го по 14-е число каждого месяца необходимо выплатить часть долга;

— 15-го числа каждого месяца долг должен быть на одну и ту же величину меньше долга
на 15-е число предыдущего месяца.

Известно, что общая сумма денег, которую нужно выплатить банку за весь срок кредитования, на 13\% больше, чем сумма, взятая в кредит. Найдите $r$.

